\documentclass[blue]{beamer}
\mode<presentation> {
  \setbeamertemplate{background canvas}[vertical shading][bottom=red!10,top=blue!10]
  \usetheme{Boadilla}
  \usefonttheme[onlysmall]{serif}
}
\usepackage{pgf,pgfarrows,pgfnodes,pgfautomata,pgfheaps,pgfshade}
\usepackage{tikz}
\usepackage{amsmath,amssymb,bm,mathtools}
\usepackage{colortbl}
\usepackage[english]{babel}
\setbeamercovered{dynamic}
\DeclareMathOperator\PV{PV}
\def\R{\mathbb{R}}
\def\Q{\mathbb{Q}}
\def\N{\mathbb{N}}
\def\E{\PV}
\newcommand{\ps}[2]{ \langle #1 , #2 \rangle }
\def\1{\mathbf{1}}
\newcommand{\nor}[1]{\left\| #1 \right\|}
\def\leq{\leqslant}
\def\geq{\geqslant}
\def\et{\quad \textrm{and} \quad}
\def\ou{\quad \textrm{or} \quad}
\def\boite{\hskip 0.5mm {\Box} \hskip 0.5mm}
\DeclareMathOperator\PPE{PPE}
\DeclareMathOperator\vNM{vNM}
\DeclareMathOperator\argmax{argmax}
\DeclareMathOperator\Prob{Prob}
\DeclareMathOperator\supp{supp}
\DeclareMathOperator\Eq{Eq}
\DeclareMathOperator\inte{int}
\DeclareMathOperator\co{co}
\DeclareMathOperator\cl{cl}
\DeclareMathOperator\Dirac{Dirac}
\DeclareMathOperator\F{F}
\DeclareMathOperator\As{As}
\DeclareMathOperator\Ir{Ir}
\DeclareMathOperator\MRS{MRS}
\DeclareMathOperator\Range{Im}
\DeclareMathOperator\DIV{DIV}
\DeclareMathOperator\A{D}
%------------------------------ theorem styles
\theoremstyle{definition}
\newtheorem*{assumption}{Assumption}
\newtheorem*{conjecture}{Conjecture}
\theoremstyle{example}
\newtheorem*{proposition}{Proposition}
\newtheorem*{remark}{Remark}
%-----------------------------------------
\title[Corporate Finance]{Theory of Corporate Finance}
\author[Victor Filipe]{V. Filipe Martins-da-Rocha}
\institute[FGV EESP]{Escola de Economia de S\~ao Paulo}
\date[April--June, 2026]{Part 1.c. Stock Market Equilibrium: About competitive price perceptions}
\subject{Economic Theory}
\begin{document}
\frame{\titlepage}
%------------------------------------------------
\begin{frame}\frametitle{Working paper}
\begin{thebibliography}{8}
  \beamertemplatearticlebibitems
    \bibitem{BisinGottardiRutta}
    Bisin, A., Gottardi. P. and Rutta, G.
    \newblock Equilibrium Corporate Finance.
    \newblock {\em NYU Working Paper} 2009
    \bibitem{BisinGottardiRutta}
    Bisin, A., Gottardi. P. and Rutta, G.
    \newblock Equilibrium Corporate Finance: Makowski meets Prescott and Townsend
    \newblock {\em EUI Working Paper} 2011
    \end{thebibliography}
\end{frame}
%------------------------------------------------
\begin{frame}\frametitle{The model: markets}
\begin{itemize}
\item One firm $J=\{j\}$ (or a continuum of identical firms)
\item One riskless financial asset $A=\{a_0\}$
    \begin{itemize}
    \item The firm can issue debt (or short-sell the financial asset $a_0$) but cannot invest: $\beta \geq 0$
    \item Investors cannot issue debt but can invest (or buy the firm's debt): $\theta^i \geq 0$
    \end{itemize}
\item Stock market
    \begin{itemize}
    \item Investors cannot short-sell the firm's shares: $\alpha^i \geq 0$
    \end{itemize}
\end{itemize}
\end{frame}
%------------------------------------------------
\begin{frame}\frametitle{The model: the firm}
\begin{itemize}
\item The production technology of the firm is defined as follows:
\[
Y^j \equiv \{(y_0,(y_s)_{s\in S}) \in \R \times \R^S \ \colon \ y_0 \leq 0 \et y_s\leq f_s(-y_0) \}
\]
where
\[
f_s :[0,\infty) \to [0,\infty)
\]
The investment $-y_0$ is denoted by $k$ representing units of commodity invested in capital at $t=0$
\end{itemize}
\end{frame}
%------------------------------------------------
\begin{frame}\frametitle{The model: the firm}
\begin{assumption}
The firm chooses the investment $k$ and the debt level $\beta$ satisfying the solvency requirement
\[
\inf_{s\in S} f_s(k) \geq \beta
\]
\end{assumption}
We denote by $\A$ the set of admissible capital budgeting and financial strategy defined by
\[
\A \equiv \{(k,\beta) \in \R^2_+ \ \colon \ \inf_{s\in S} f_s(k) \geq \beta \}
\]
\end{frame}
%------------------------------------------------
\begin{frame}\frametitle{Budget restrictions for agents}
Each investor~$i$ takes as given
\begin{itemize}
\item the firm's capital budgeting and financial strategy $(k,\beta)$
\item the security price $q$ and the stock (equity) price $E$
\end{itemize}
and chooses
\begin{itemize}
\item a consumption plan $c=(c_0,c_1) \in \R_+ \times \R_+^S$
\item an investment $\theta^i\geq 0$ and a share $\alpha^i \geq 0$ satisfying
\begin{itemize}
\item the budget restriction at $t=0$
\[
c^i_0 + q \theta^i + E \alpha^i \leq \omega^i_0 + [E+q\beta - k] \xi^i
\]
\item the budget restriction at $t=1$ for each state $s$
\[
c^i_s \leq \omega^i_s + \theta^i + [f_s(k) - \beta] \alpha^i
\]
\end{itemize}
\item We denote by $B^i(q,E,k,\beta)$ the set of all actions $(c,\theta,\alpha)$ satisfying the budget restrictions
\end{itemize}
\end{frame}
%------------------------------------------------
\begin{frame}\frametitle{Competitive equilibrium given the firm's decisions}
Fix a firm's decision $(k,\beta) \in \A$, a family
\[
\{(k,\beta),(q,E),(c^i,\theta^i,\alpha^i)_{i\in I}\}
\]
is a competitive equilibrium if
\begin{itemize}
\item every agent~$i$'s action $(c^i,\theta^i,\alpha^i)$ is optimal, i.e.,
\[
(c^i,\theta^i,\alpha^i) \in \argmax \{U^i(c) \ \colon \ (c,\theta,\alpha) \in B^i(q,E,\beta,k) \}
\]
\item financial markets clear, i.e.,
\[
\sum_{i\in I} \theta^i = \beta \et \sum_{i\in I} \alpha^i = 1
\]
\item consumption markets clear, i.e.,
\[
k + \sum_{i\in I} c^i_0 = \omega_0 \et \sum_{i\in I} c^i_1 = \omega_1 + f(k)
\]
\end{itemize}
\end{frame}
%------------------------------------------------
\begin{frame}\frametitle{Notations}
Fix a consumption plan  $c^i=(c^i_0,c^i_1)$ satisfying
\[
c^i_0 >0 \et c^i_s >0, \ \forall s\in S
\]
Given a random variable $w=(w_s)_{s\in S} \in \R^S$, we let
\[
\E^i[w] = \sum_{s\in S} \MRS^i_{s,0}(c^i) w_s
\]
where
\[
\MRS^i_{s,0}(c^i) \equiv \beta^i \nu^i_s \frac{\partial u^i(c^i_s)}{\partial u^i(c^i_0)}
\]
\end{frame}
%------------------------------------------------
\begin{frame}\frametitle{Competitive equilibrium: necessary conditions}
\begin{proposition}
Assume there exists a competitive equilibrium
\[
\{(k,\beta),(q,E),(c^i,\theta^i,\alpha^i)_{i\in I}\} \quad \textrm{with} \quad \beta >0
\]
Then
\[
q = \max_{i\in I} \sum_{s\in S} \MRS^i_{s,0}(c^i) = \max_{i\in I} \E^i \left[ \1_S \right]
\]
and
\[
E = \max_{i\in I} \sum_{s\in S} \MRS^i_{s,0}(c^i) \{f_s(k)-\beta\} = \max_{i\in I} \E^i  \left[  f_1(k)- \beta \1_S \right]
\]
\end{proposition}
\end{frame}
%------------------------------------------------
\begin{frame}\frametitle{Investors' price perception: definition}
Fix a competitive stock market equilibrium
\[
\{(k,\beta),(q,E),(c^i,\theta^i,\alpha^i)_{i\in I}\}
\]
If the firm production and financing choice were different, say $(k',\beta')$ then its cash flow would be $(f_s(k') - \beta')_{s\in S}$ and consequently, the stock price should be different
\end{frame}
%------------------------------------------------
\begin{frame}\frametitle{Investors' price perception: Makowski criteriom}
We assume that each investor
    \begin{itemize}
    \item is aware that the stock price is a function of $(k',\beta')$
    \item believes (conjectures) that the value of the equity price map $$(k',\beta') \mapsto \widetilde{E}(k',\beta')$$ is the highest marginal valuation across all investors of the firm's cash flow
associated to $(k',\beta')$, i.e.,
        \[
        \widetilde{E}(k',\beta') = \max_{i \in I} \E^i \left[ f_1(k')-\beta' \1_S \right]
        \]
    \end{itemize}
\end{frame}
%------------------------------------------------
\begin{frame}\frametitle{Makowski}
\begin{thebibliography}{8}
  \beamertemplatearticlebibitems
    \bibitem{Makowski80}
    Makowski, L.
    \newblock A characterization of perfectly competitive economies with production
    \newblock {\em Journal of Economic Theory} 22, 1980
    \bibitem{Makowski83a}
    Makowski, L.
    \newblock Competition and unanimity revisited
    \newblock {\em American Economic Review} 73, 1983
    \bibitem{Makowski83b}
    Makowski, L.
    \newblock Competitive stock markets
    \newblock {\em Review of Economic Studies} 50, 1983
    \end{thebibliography}
\end{frame}
%------------------------------------------------
\begin{frame}\frametitle{Investors' price perception: remarks}
\begin{itemize}
\item At equilibrium investors' price perception is consistent with the equilibrium stock price, i.e.,
\[
E = \widetilde{E}(k,\beta)
\]
\item The equity price map $\widetilde{E}$ takes into account the fact that the firm's choice $(k,\beta)$ modifies its cash flow, but it is assumed that investors' marginal rates of substitution $\MRS^i(c^i)$ used to determine the market valuation are independent of such choice of $(k,\beta)$
\end{itemize}
\begin{block}{Justification}
Continuum of identical firms: one firm's decision has a negligible impact on market prices since agents cannot short-sell stocks
\end{block}
\end{frame}
%------------------------------------------------
\begin{frame}\frametitle{Investors' price perception: remarks}
\begin{itemize}
\item Agent~$i_0$ valuates a firm's cash flow thinking of the benefits he may obtain selling the firm in the market
\item Therefore, agent~$i_0$ uses the MRS of those investors willing to pay the most
\end{itemize}
\begin{block}{Problem}
Every agent should know (or be able to compute) the MRS of each investor
\end{block}
\end{frame}
%------------------------------------------------
\begin{frame}\frametitle{Strong unanimity: result}
For a firm's decision $(k',\beta')$, we denote by $B^i(q,E',k',\beta')$ the budget set associated to the security price $q$ and the equity price $E' = \widetilde{E}(k',\beta')$
\begin{theorem}[Strong unanimity]
Fix a competitive stock market equilibrium
\[
\{(k,\beta),(q,E),(c^i,\theta^i,\alpha^i)_{i\in I}\}
\]
where the firm has chosen $(k,\beta)$ to maximize its value defined by
\[
\forall (k',\beta') \in \A,\quad V(k',\beta') \equiv \widetilde{E}(k',\beta') + q \beta' - k'
\]
For each investor~$i$, we have
\[
U^i(c) = \max \{ U^i(c') \ \colon \ (c',\theta',\alpha') \in B^i(q,E',k',\beta') \et (k',\beta') \in \A \}
\]
\end{theorem}
In particular,
\[
(k,\beta) \in \argmax \{ U^i(c') \ \colon \ (c',\theta',\alpha') \in B^i(q,E',k',\beta') \et (k',\beta') \in \A \}
\]
\end{frame}
%------------------------------------------------
\begin{frame}\frametitle{Strong unanimity: comments}
\begin{itemize}
\item \alert{Unanimity}: Every agent agrees that the manager should maximize the function $(k,\beta) \mapsto \widetilde{E}(k,\beta) + q\beta -k$
\item \alert{Strong}: When agent~$i$ considers possible deviations $(k',\beta')$ of the manager, he also considers possible deviations of his own choices $(c',\theta',\alpha')$
\end{itemize}
\end{frame}
%====================================================================
\section{Constrained Pareto Efficiency}
%====================================================================
\begin{frame}\frametitle{Constrained Pareto efficiency: admissible allocations}
\begin{definition}
An allocation $(c^i)_{i\in I}$ of consumption plans is said \alert{admissible} if
\begin{itemize}
\item it is feasible: there exists an investment level $k\geq 0$ such that
\[
k + \sum_{i\in I} c^i_0 \leq \omega_0 \et \sum_{i\in I} c^i_1 \leq f(k) +  \sum_{i\in I} \omega^i_1
\]
\item it is attainable: there exists $\beta \in [0, \inf_{s\in S} f_s(k)]$ and, for each~$i$, $(\theta^i,\alpha^i) \in \R_+^2$ satisfying
    \[
    c^i_s = \omega^i_s + [f_s(k) - \beta] \alpha^i + \theta^i, \quad \forall s\in S
    \]
\end{itemize}
\end{definition}
\bigskip
The financial portfolios $(\theta^i, \alpha^i)$ are \emph{not} required to clear in aggregate
\end{frame}
%------------------------------------------------
\begin{frame}\frametitle{Weak Pareto efficiency}
\begin{definition}[Weak Pareto efficiency]
An admissible allocation $(c^i)_{i\in I}$ is \alert{weakly constrained Pareto efficient} if there is no other admissible allocation $(\widetilde{c}^i)_{i\in I}$ with
\[
U^i(\widetilde{c}^i) > U^i(c^i), \quad \text{for all $i\in I$}
\]
\end{definition}
\bigskip
The (strong) notion would forbid $U^i(\widetilde{c}^i)\geq U^i(c^i)$ for all $i$ with strict inequality for some $i$
\bigskip

Under continuity and strict monotonicity of $U^i$, weak and strong Pareto efficiency \alert{coincide} \\
\hfill\small(Mas-Colell, Whinston and Green, 1995, Prop.~16.D.2)
\end{frame}
%------------------------------------------------
\begin{frame}\frametitle{Weak welfare theorem}
\begin{theorem}[Weak welfare theorem]
Every competitive equilibrium allocation is weakly constrained Pareto efficient.
\end{theorem}
\bigskip

\textbf{Proof idea.} Suppose by contradiction that $(\widetilde{c}^i)_{i \in I}$ admissible Pareto-dominates the equilibrium in the strict sense.

For each~$i$, define the \emph{wealth gap}
\[
G^i \equiv \omega^i_0 + V(\widetilde{k}, \widetilde{\beta}) \xi^i - \widetilde{c}^i_0 - q \widetilde{\theta}^i - \widetilde{E}(\widetilde{k}, \widetilde{\beta}) \widetilde{\alpha}^i
\]

By strong unanimity, $G^i < 0$ for every~$i$. But aggregating $G^i$ and using market clearing yields $\sum_i G^i \geq 0$. Contradiction.
\end{frame}
%====================================================================
\section{Modigliani--Miller Revisited}
%====================================================================
\begin{frame}\frametitle{Equity holders and bond holders}
Fix a competitive equilibrium $\{(k,\beta),(q,E),(c^i,\alpha^i,\theta^i)_{i\in I}\}$
\begin{definition}[Marginal holders]
For any $\beta' \in [0, \overline{\beta}]$ where $\overline{\beta} \coloneqq \inf_{s \in S} f_s(k)$:
\[
I^e(\beta') \coloneqq \{i \in I \colon \widetilde{E}(k, \beta') = \PV^i[f_1(k) - \beta'\1_S]\}
\]
\[
I^b \coloneqq \{i \in I \colon q = \PV^i[\1_S]\}
\]
\end{definition}
\bigskip
\begin{itemize}
\item $i \in I^e(\beta')$: \emph{equity holder} at the conjectured price $\widetilde{E}(k,\beta')$
\item $i \in I^b$: \emph{bond holder}
\end{itemize}
\bigskip

By max-valuation pricing, $I^e(\beta')$ and $I^b$ are non-empty
\end{frame}
%------------------------------------------------
\begin{frame}\frametitle{Convexity of the value map}
\begin{lemma}[Convexity]
For each $k \geq 0$, the value map
\[
\beta' \mapsto V(k, \beta') = \widetilde{E}(k, \beta') + q\beta' - k
\]
is convex piecewise linear on $[0, \overline{\beta}]$.
\end{lemma}
\bigskip

\textbf{Why?} For each $i \in I$, the function
\[
g_i(\beta') \coloneqq \PV^i[f_1(k) - \beta' \1_S] = \PV^i[f_1(k)] - \beta' \PV^i[\1_S]
\]
is \emph{affine} in $\beta'$
\bigskip

$\widetilde{E}(k, \beta') = \max_{i \in I} g_i(\beta')$ is the max of affine functions: \alert{convex piecewise linear}
\bigskip

Adding the linear term $q\beta' - k$ preserves convexity
\end{frame}
%------------------------------------------------
\begin{frame}\frametitle{Monotonicity of the value map}
By max-valuation pricing, $\PV^i[\1_S] \leq q$ for every~$i$
\bigskip

Each function $g_i(\beta')$ has slope $g_i'(\beta') = -\PV^i[\1_S] \in [-q, 0]$
\bigskip

The slope of $\widetilde{E}(k, \cdot) = \max_i g_i$ lies in $[-q, 0]$ wherever it is differentiable
\bigskip

Adding $q$ for the linear term:
\[
\text{slope of $V(k, \cdot)$ lies in $[0, q]$}
\]
\bigskip

\alert{$V(k, \cdot)$ is non-decreasing on $[0, \overline{\beta}]$}
\end{frame}
%------------------------------------------------
\begin{frame}\frametitle{The Modigliani--Miller dichotomy}
\begin{proposition}[MM dichotomy]
For each $k \geq 0$, exactly one of the following holds:
\bigskip

\begin{itemize}
\item[\textup{(MM)}] $V(k, \cdot)$ is constant on $[0, \overline{\beta}]$. Every $\beta$ is value-maximizing. Capital structure is fully indeterminate.
\bigskip

\item[\textup{($\neg$MM)}] $V(k, \cdot)$ is not constant. The unique value-maximizer is $\overline{\beta}$.
\end{itemize}
\end{proposition}
\bigskip

\textbf{Proof:} A convex function attaining its maximum at an interior point is constant on the whole interval. Combined with monotonicity ($V$ non-decreasing), the only non-constant case has max at the right endpoint $\overline{\beta}$.
\end{frame}
%------------------------------------------------
\begin{frame}\frametitle{Three possible shapes of $V(k, \cdot)$}
\begin{center}
\begin{tikzpicture}[scale=0.85, font=\small]
  % Panel A: V constant
  \begin{scope}
    \draw[->, gray] (-0.2, 0) -- (3.4, 0) node[below right, black] {\scriptsize $\beta$};
    \draw[->, gray] (0, -0.2) -- (0, 2.8);
    \draw[ultra thick, blue] (0, 1.5) -- (3, 1.5);
    \node[below] at (3, 0) {\scriptsize $\overline{\beta}$};
    \node[below] at (1.5, -0.5) {\textbf{(MM)}};
    \node[below] at (1.5, -1.0) {\scriptsize constant};
  \end{scope}
  % Panel B: V strictly increasing
  \begin{scope}[xshift=4.7cm]
    \draw[->, gray] (-0.2, 0) -- (3.4, 0) node[below right, black] {\scriptsize $\beta$};
    \draw[->, gray] (0, -0.2) -- (0, 2.8);
    \draw[ultra thick, blue] (0, 0.5) -- (3, 2.5);
    \node[below] at (3, 0) {\scriptsize $\overline{\beta}$};
    \node[below] at (1.5, -0.5) {\textbf{($\neg$MM)}};
    \node[below] at (1.5, -1.0) {\scriptsize strictly increasing};
  \end{scope}
  % Panel C: flat then increasing
  \begin{scope}[xshift=9.4cm]
    \draw[->, gray] (-0.2, 0) -- (3.4, 0) node[below right, black] {\scriptsize $\beta$};
    \draw[->, gray] (0, -0.2) -- (0, 2.8);
    \draw[ultra thick, blue] (0, 0.7) -- (1.2, 0.7) -- (3, 2.5);
    \draw[dotted, gray] (1.2, 0.7) -- (1.2, 0);
    \node[below] at (1.2, 0) {\scriptsize $\beta^\ast$};
    \node[below right] at (3, 0) {\scriptsize $\overline{\beta}$};
    \node[below] at (1.5, -0.5) {\textbf{($\neg$MM)}};
    \node[below] at (1.5, -1.0) {\scriptsize flat then increasing};
  \end{scope}
\end{tikzpicture}
\end{center}
\bigskip

The opposite shape ``strictly increasing then constant'' is \alert{ruled out by convexity}: it would require a concave kink
\bigskip

In all $(\neg$MM$)$ cases, the firm's unique optimum is $\overline{\beta}$ \\(possibly with a flat segment $[0, \beta^\ast]$ that does \emph{not} contain the optimum)
\end{frame}
%------------------------------------------------
\begin{frame}\frametitle{Equity holders are bond holders in the MM regime}
\begin{proposition}[Equity holders are bond holders]
If $V(k, \cdot)$ is constant on $[0, \overline{\beta}]$ \textup{(}case~\textup{(MM)} of the dichotomy\textup{)}, then
\[
I^e(\beta') \subset I^b, \quad \text{for every $\beta' \in [0, \overline{\beta})$}
\]
\end{proposition}
\bigskip

The proof is \alert{direct}, without computing directional derivatives. It uses only:
\begin{enumerate}
\item Definition of $I^e$ as the active set of $\widetilde{E}(k, \cdot)$
\item Makowski maximum dominates individual valuations
\item Constancy of $V$ on the no-default interval
\item Max-valuation pricing of the bond
\end{enumerate}
\end{frame}
%------------------------------------------------
\begin{frame}\frametitle{Proof}
Fix $\beta' \in [0, \overline{\beta})$ and $i \in I^e(\beta')$. By definition:
\[
\widetilde{E}(k, \beta') = \PV^i[f_1(k)] - \beta' \PV^i[\1_S]
\]
For any $\beta'' \in (\beta', \overline{\beta}]$, Makowski maximum dominates:
\[
\widetilde{E}(k, \beta'') \geq \PV^i[f_1(k)] - \beta'' \PV^i[\1_S]
\]
Subtracting:
\[
\widetilde{E}(k, \beta'') - \widetilde{E}(k, \beta') \geq (\beta' - \beta'') \PV^i[\1_S]
\]
By constancy of $V$: \quad $\widetilde{E}(k, \beta'') - \widetilde{E}(k, \beta') = q(\beta' - \beta'')$
\bigskip

Combining: \quad $q(\beta' - \beta'') \geq (\beta' - \beta'') \PV^i[\1_S]$
\bigskip

Since $\beta' - \beta'' < 0$, dividing reverses the inequality: $q \leq \PV^i[\1_S]$
\bigskip

Combined with max-valuation $\PV^i[\1_S] \leq q$: \quad $\PV^i[\1_S] = q$, i.e., $i \in I^b$ \alert{$\boite$}
\end{frame}
%------------------------------------------------
\begin{frame}\frametitle{The two regimes: economic content}
\textbf{(MM) regime:}
\begin{itemize}
\item Capital structure fully indeterminate
\item Every active equity holder is a bond holder
\item Market values bond and equity at consistent rates: $\PV^i[\1_S] = q$ for all $i \in I^e(\beta')$ at every $\beta' \in [0, \overline{\beta})$
\end{itemize}
\bigskip

\textbf{$(\neg$MM$)$ regime:}
\begin{itemize}
\item Capital structure determinate: firm pushes debt to $\overline{\beta}$
\item No-default constraint binds at equilibrium
\item There exists $\beta' < \overline{\beta}$ where some equity holder strictly under-values the bond: $\PV^i[\1_S] < q$ for some $i \in I^e(\beta')$
\end{itemize}
\bigskip

The intermediate case ``firm indifferent on a strict sub-interval that does not include the endpoints'' is \alert{impossible} by convexity
\end{frame}

% %------------------------------------------------
% \begin{frame}\frametitle{Partial indeterminacy: a family of equilibria}
% A different notion of partial indeterminacy appears in \alert{Bisin--Gottardi--Ruta (2014, Fig.~1)}:
% \bigskip

% A continuum of equilibria parameterized by $B^\ast \in [B_{\min}, \overline{\beta}]$, all sharing
% \begin{itemize}
% \item the same consumption allocations
% \item the same individual valuations $\PV^i$
% \item the same value function $V(k, \cdot)$, constant on $[0, \overline{\beta}]$
% \end{itemize}
% \bigskip

% Each equilibrium in the family is in the (MM) regime. The indeterminacy is \alert{across} the family, not within a single equilibrium
% \bigskip

% \textbf{Construction (strong MM, upward).} From the equilibrium with $B_{\min}$, adjust
% \[
% \theta'^i = \theta^i + (B^\ast - B_{\min}) \alpha^i, \quad \alpha'^i = \alpha^i
% \]
% to obtain a new equilibrium with $B^\ast$, same consumption and prices

% \bigskip
% In the BGR parameterization, $B_{\min} = 0.1828$
% \end{frame}
% %====================================================================
% \section{Intermediated Short Sales}
% %====================================================================
% \begin{frame}\frametitle{Intermediated short sales}
% \begin{itemize}
% \item Consider now the case where consumers can sell short the firm's equity
% \item A short sale is a loan contract with a promise to repay an amount equal to the future value of equity (here it coincides with dividend)
% \item We introduce financial intermediaries who can issue claims corresponding to both short and long positions on the firm's equity
% \item Intermediaries bear no cost to issue claims, but face the possibility of default on the short positions they issue
% \item To simplify the analysis, we consider a reduced form version of the model:
%     \begin{itemize}
%     \item the default rate on short positions is exogenously given and equal to $\delta>0$ in every state, for all consumers
%     \item the analysis should be extended to the general case where default rates are endogenously chosen by consumers
%     \end{itemize}
% \end{itemize}
% \end{frame}
% %------------------------------------------------
% \begin{frame}\frametitle{Intermediated short sales}
% \begin{itemize}
% \item Intermediating $m$ units of the derivative the intermediary is repaid only a fraction $(1-\delta)$ of the amount due on each short position issued
% \item To insure his own solvency, the intermediary is required to hold an appropriate amount of equity, as a form of collateral, to cover the shortfalls in the revenue
% \item The amount $\gamma(m)$ of equity of the firm that the intermediary must hold should satisfy
% \[
% m \leq m(1-\delta) + \gamma(m)
% \]
% \item This collateral cost can be covered by different (competitive) prices for long and short positions
%     \begin{itemize}
%     \item $q^+$ is the price at which long positions (on the derivative) are traded
%     \item $q^-$ is the price at which short positions (on the derivative) are traded
%     \end{itemize}
% \end{itemize}
% \end{frame}
% %------------------------------------------------
% \begin{frame}\frametitle{Intermediated short sales: the intermediary problem}
% The intermediary chooses
% \begin{itemize}
% \item the amount $m \geq 0$ of long and short positions
% \item the amount $\gamma$ of equity held as a hedge
% \end{itemize}
% so as to maximize his total revenue at $t=0$
% \[
% \max \{ (q^+ - q^-)m - q\gamma \ \colon\ m\geq 0 \et m \leq m(1-\delta) + \gamma \}
% \]
% The intermediary chooses
% \[
% \gamma(m) = m\delta
% \]
% and $m>0$ arbitrary only if $q=(q^+ - q^-)/2$
% \begin{remark}
% The intermediary technology displays constant returns to scale leading to zero profits
% \end{remark}
% \end{frame}
% %------------------------------------------------
% \begin{frame}\frametitle{Intermediated short sales: consumer's budget restriction}
% Let
% \begin{itemize}
% \item $\phi^i_{+} \geq 0$ denote consumer~$i$'s holdings of long positions in the derivative
% \item $\phi^i_{-} \geq 0$ denote consumer~$i$'s holdings of long positions in the derivative
% \end{itemize}
% The consumer budget constraints are then
% \[
% c^i_0 + q \theta^i + q^{+} \phi^i_{+} + E \alpha^i  \leq \omega^i_0 + [E+q\beta - k] \xi^i + q_{-} \phi^i_{-}
% \]
% and for each~$s$ at $t=1$
% \[
% c^i_s + [f_s(k) - \beta] \phi^i_{-}(1-\delta) \leq \omega^i_s + [f_s(k) - \beta](\alpha^i + \phi^i_{+}) + \theta^i
% \]
% \end{frame}
% %------------------------------------------------
% \begin{frame}\frametitle{Intermediated short sales: market clearing}
% \begin{itemize}
% \item The equity market clearing is now
% \[
% \gamma + \sum_{i\in I} \theta^i = 1
% \]
% \item The derivative market clearing is
% \[
% \sum_{i\in I} \phi^i_{+} = \sum_{i\in I} \phi^i_{-} = m
% \]
% where $\gamma = \delta m$
% \end{itemize}
% \end{frame}
% %------------------------------------------------
% \begin{frame}\frametitle{Intermediated short sales: price conjectures}
% \begin{itemize}
% \item Intermediaries, in addition to consumers, may demand equity in the market
% \item The price conjecture $\tilde{E}(k,\beta)$ should properly be adjusted
% \end{itemize}
% We pose
% \[
% \tilde{E}(k,\beta) = \max \left\{\tilde{E}^c(k,\beta),\underbrace{\frac{1}{\beta}[\tilde{q}^+(k,\beta) - \tilde{q}^{-}(k,\beta)]}_{\textrm{value of intermediation}} \right\}
% \]
% where
% \[
% \tilde{E}^c(k,\beta) = \max_{i\in I} \PV^i[f_1(k) - \beta \1_S] = \max_{i\in I} \sum_{s\in S} \MRS^i_s (f_s(k) - \beta)
% \]
% and
% \begin{eqnarray*}
% \tilde{q}^{+}(k,\beta) & = & \max_{i\in I} \PV^i[f_1(k) - \beta \1_S] \\
% \tilde{q}^{-}(k,\beta)  & = &  (1-\delta) \min_{i\in I} \PV^i[f_1(k) - \beta \1_S]
% \end{eqnarray*}
% \end{frame}
% %------------------------------------------------
% \begin{frame}\frametitle{Intermediated short sales: equilibrium}
% Two situations may occur at equilibrium with intermediation ($m>0$)
% \[
% q > q^+ \ou q= q^+
% \]
% \begin{block}{$q > q^+$}
% \begin{itemize}
% \item Equity sells at a premium over the long positions on the derivative
% \item This is because of its additional value as input in the intermediation technology
% \item All amount of equity is purchased by the intermediary
% \item This is in turn implies that $q=(q^+ - q^{-})/\delta$
% \item The additional cost of equity is compensated by the presence of a sufficiently high spread $q^+ - q^{-}$
% \end{itemize}
% \end{block}
% \end{frame}
% %------------------------------------------------
% \begin{frame}\frametitle{Intermediated short sales: equilibrium}
% Two situations may occur at equilibrium with intermediation ($m>0$)
% \[
% q > q^+ \ou q= q^+
% \]
% \begin{block}{$q = q^+$}
% \begin{itemize}
% \item Equity and the long positions on the derivative are traded at the same price
% \item Consumers are indifferent between buying long positions in equity and the derivative
% \item Some of outstanding amount of equity is held by consumers
% \end{itemize}
% \end{block}
% \end{frame}
% %------------------------------------------------
% \begin{frame}\frametitle{Intermediated short sales: unanimity and efficiency}
% \begin{proposition}
% At a competitive equilibrium of an economy with intermediated short-sales, equity holders unanimously support the production and financial decisions of the firms.
% Moreover, the equilibrium allocation is constrained Pareto optimal
% \end{proposition}
% \end{frame}
% %------------------------------------------------
% \begin{frame}\frametitle{Bisin, Gottardi and Rutta}
% \begin{itemize}
% \item Risky debt
% \item Asymmetric information
% \end{itemize}
% \end{frame}

%------------------------------------------------
\begin{frame}\frametitle{Literature}
{\small
\begin{thebibliography}{8}
  \beamertemplatearticlebibitems
    \bibitem{AcharyaBisin09}
    Acharya, V.V. and A. Bisin
    \newblock Managerial hedging, equity ownership, and firm value
    \newblock {\em RAND Journal of Economics} 40, 2009
    \bibitem{AllenGale91}
    Allen, F. and D. Gale
    \newblock Arbitrage, short sales, and financial innovation
    \newblock {\em Econometrica} 59, 1991
    \bibitem{BisinGottardiRampini08}
    Bisin, A., P. Gottardi, and A. Rampini
    \newblock Managerial hedging and portfolio monitoring
    \newblock {\em Journal of the European Economic Association} 6, 2008
    \bibitem{BisinGuaitoli04}
    Bisin, A. and D. Guaitoli
    \newblock Moral Hazard and Non-exclusive Contracts
    \newblock {\em RAND Journal of Economics} 35, 2004
    \bibitem{Carceles-PovedaCoenPirani09}
    Carceles-Poveda, E. and D. Coen-Pirani
    \newblock Shareholders' unanimity with incomplete markets
    \newblock {\em International Economic Review} 50, 2009
    \end{thebibliography}
}
\end{frame}
\end{document}